% Generated draft: requires mathematical and source review.
\documentclass[11pt]{article}
\usepackage[margin=25mm]{geometry}
\usepackage{fontspec}
\setmainfont{Latin Modern Roman}
\usepackage{amsmath,amssymb,mathtools,unicode-math}
\setmathfont{Latin Modern Math}
\usepackage{enumitem,needspace,xurl,hyperref}
\setlength{\parindent}{0pt}
\setlength{\parskip}{5pt}
\setlength{\emergencystretch}{3em}
\newcommand{\N}{\mathbb{N}}
\newcommand{\Z}{\mathbb{Z}}
\newcommand{\Q}{\mathbb{Q}}
\newcommand{\R}{\mathbb{R}}
\newcommand{\C}{\mathbb{C}}
\newcommand{\F}{\mathbb{F}}
\newcommand{\A}{\mathbb{A}}
\newcommand{\PP}{\mathbb{P}}
\newcommand{\E}{\mathbb{E}}
\newcommand{\eps}{\varepsilon}
\newcommand{\dd}{\,\mathrm{d}}
\begin{document}
\section*{UnsolvedMath ID: 700770}
\textbf{Source problem code:} catalogue entry
\subsection{Definitions and mathematical statement}
For an algebraic group \(G\) defined over a field \(k\), write \(G(k)\) for its abstract group of \(k\)-rational points. A field \(k\) is called \emph{algebraic} if every element of \(k\) is algebraic over its prime subfield. An algebraic group is directly undecomposable, in the usual algebraic-group sense, if it is not isomorphic to a direct product of two nontrivial algebraic groups. A trivial centre means that the centre consists only of the identity. The source does not specify the category in which direct decomposability and the centre are to be interpreted; this qualification remains unresolved.

The source calls an abstract group isomorphism between groups of points \emph{semialgebraic} if it is a composition of an ``isomorphism of translation of the field of definition'' and a rational morphism. The natural interpretation of the field-translation step is transport by a field isomorphism \(\sigma:k\to\ell\). Writing \({}^{\sigma}G\) for the algebraic group obtained by transporting the defining data of \(G\) along \(\sigma\), this interpretation gives a factorization
\[
G(k)\xrightarrow{\sigma_*}({}^{\sigma}G)(\ell)\xrightarrow{f_*}H(\ell),
\qquad \varphi=f_*\circ\sigma_*,
\]
where \(f\) is the rational morphism intended by the source and \(f_*\) must be defined on all points needed for this factorization. The source does not explain the quoted field-translation phrase or specify whether its rational morphism is a rational map, an everywhere-defined morphism, or an isomorphism of algebraic groups. This factorization is therefore an interpretation of the supplied definition requiring review, rather than a resolution of those distinctions.

\begin{enumerate}
\item[(a)] Let \(k\) be an algebraically closed field, and let \(G\) and \(H\) be directly undecomposable algebraic groups over \(k\), both with trivial centres. Does the existence of an abstract group isomorphism \(G(k)\cong H(k)\) imply the existence of a semialgebraic isomorphism between these groups of points? Equivalently, with the preceding qualifications,
\[
\operatorname{Iso}_{\mathrm{grp}}(G(k),H(k))\ne\varnothing
\quad\Longrightarrow\quad
\exists\psi\in\operatorname{Iso}_{\mathrm{grp}}(G(k),H(k))
\text{ such that }\psi\text{ is semialgebraic}.
\]
The assertion concerns the existence of a semialgebraic isomorphism; it does not require every given isomorphism to be semialgebraic.
\item[(b)] Is every abstract group isomorphism between the groups of points of directly undecomposable algebraic groups with trivial centres, defined over algebraic fields, semialgebraic? If separate fields of definition are allowed, the quantified question is whether, for every pair of algebraic fields \(k,\ell\), every such \(G/k\) and \(H/\ell\), and every
\[
\varphi\in\operatorname{Iso}_{\mathrm{grp}}(G(k),H(\ell)),
\]
\(\varphi\) is semialgebraic. If a common field of definition is intended, the question instead restricts this formulation to \(k=\ell\). The supplied wording does not explicitly distinguish these two scopes. Algebraic fields in this part are not assumed to be algebraically closed.
\end{enumerate}
\subsection{Short English statement}
Part (a) asks whether abstractly isomorphic point groups of directly undecomposable, centreless algebraic groups over an algebraically closed field always admit some semialgebraic isomorphism. Part (b) asks whether every such abstract isomorphism over algebraic fields is semialgebraic. The precise meaning of the field-translation step, the rational morphism, and the field relationship in part (b) requires clarification.
\subsection{Sources}
\begin{itemize}
\item Local UnsolvedMath registry, supplied complete source record: permanent ID \(700770\), exact source problem code \texttt{catalogue entry}, title \emph{Kourovka Notebook Problem 11.78}; upstream revision \texttt{d371ee917ac9647fbb06cd6fc529deec6f3dda82}. The record was created and updated on 2026-10-16 and labels the problem \emph{open}; this is supplied registry information, not an independently verified mathematical or literature assessment.
\item Attributed in the supplied record to K. N. Ponomar\"ev, proposed in 1990. Supplied reference: \emph{The Kourovka Notebook (Unsolved Problems in Group Theory)}, 11th issue (1990), Problem 11.78; also identified with the 21st issue edition (2026), \texttt{arXiv:1401.0300v48}. The supplied collection description names E. I. Khukhro and V. D. Mazurov as editors. These bibliographic details and the problem wording have not been independently checked against the cited editions.
\end{itemize}
\end{document}
